\section{Conclusion}
\label{sec:conclusion}

In 17 French departments over 26 summers, extreme heat in July and August is associated with a wider yield advantage of irrigated over rainfed grain maize: \QonePct\,\% per degree-day above 32\,$^{\circ}$C, with a test whose size was checked on the same data. Rainfed yields are about \QoneRatio{} times more sensitive to heat than irrigated yields, which are themselves not immune. The association met a decision rule fixed before the replication years were seen, kept its sign on a decade the exploratory analysis had not used, with an estimate about half the 2010--2025 one, and did not change with an older map of maize. It depends on measuring heat where maize is grown, and in the earlier decade it does not hold with department-wide weather. It is an association, not the causal effect of irrigation, and it says nothing direct about a warmer climate.

The other questions gave null results. For wheat over 1959--2025, heat at the end of the cycle is associated with lower yields, but this cannot be distinguished from zero, and no change in sensitivity over time is detectable. Forecasts of wheat yields from weather and satellite data, as pre-specified, are not significantly better than the historical mean of each department.

The most general finding is methodological. In panels where a weather shock hits many units in the same year, standard errors clustered by unit or with a spatial cut-off can reject a true null in a large share of cases: \UsualSizeMin\,\% to \UsualSizeMax\,\% here. A placebo that keeps the real geography of the shock and permutes years reveals the problem in minutes. The code for this check is released with the monograph, so that it can be run on other data before their conclusions are drawn.
